% Content of §4 Experimental Setup: section header is in the master file.

\subsection{Experimental setup}

\noindent\textbf{Evaluation protocols.} The primary analysis is the CQ-125 grounding-effect study in \S\ref{sec:gisr}: no-grounding (A) versus grounding-only (B) across 11 LLM families, with $N{=}3$ paired samples and question-level majority-vote McNemar tests. A focused \texttt{gemini-3.5-flash} diagnostic adds grounding plus minimal-modification rules (C) and uses a balanced A/B/C design with $N{=}5$. Table~\ref{tab:protocol-inventory} separates this inferential panel from supporting framework, deployment, and harness-readiness results.

\begin{table}[t]
\centering
\footnotesize
\caption{Evaluation-protocol inventory. MV denotes question-level majority vote before paired McNemar testing. A is no-grounding, B is grounding-only, and C is grounding plus the minimal-modification prompt profile.}
\label{tab:protocol-inventory}
{\setlength{\tabcolsep}{4pt}
\begin{tabular}{p{0.23\linewidth}p{0.25\linewidth}p{0.20\linewidth}p{0.23\linewidth}}
\toprule
Experiment & Dataset / model scope & Conditions and samples & Role in the manuscript \\
\midrule
Cross-family grounding panel & CQ-125; 11 LLM families & A/B, $N{=}3$, MV McNemar & Primary evidence for subset-conditional and family-conditional grounding effects. \\
Focused Gemini 3.5 diagnostic & CQ-125; \texttt{gemini-3.5-flash} & A/B/C, balanced $N{=}5$, MV McNemar & Diagnostic check of the negative Spatial tail and mini-mod candidate patch. \\
Framework context & GIS, BIRD, cross-lingual, and DIN-SQL tracks & Original paired protocols; BIRD in single-pass mode & Supporting system context, not the primary grounding-effect claim. \\
Gemma4 fixed-host sweep & CQ-125; five local Gemma4 scales on host228 & Single-host frozen summaries; best-observed full rows where multiple full runs existed & Deployment accuracy--runtime probe only. \\
FloodSQL-Bench smoke audit & Local PostGIS port; \texttt{gemini-2.5-flash} & Stratified 150q, $N{=}1$ & External harness-readiness audit only. \\
\bottomrule
\end{tabular}
}
\end{table}

Framework-level results in \S\ref{sec:baseline-results} use a separate protocol and provide context only. Early full-pipeline runs used an ADK agent loop that invoked grounding, execution, and retry tools over multiple turns, costing $\sim$32$\times$ baseline tokens and sometimes entering unbounded tool-call cycles. We therefore added a deterministic \emph{single-pass} pipeline: build grounding locally, make one SQL-generation LLM call, postprocess and execute, then retry with a corrective LLM up to $T=2$ if execution fails. Single-pass reduced BIRD token cost to $7.9\times$ baseline and improved design-set EX; BIRD results use this mode, while the earlier GIS framework headline remains separated from the cross-family grounding-effect analysis.

We also include a Gemma4 size sweep on the same local Ollama node (host228), covering \texttt{e2b}, \texttt{e4b}, \texttt{12b}, \texttt{26b}, and \texttt{31b}. It is a single-host engineering probe of accuracy--runtime trade-offs and is excluded from the $N{=}3$ majority-vote McNemar panel.

\medskip
\noindent\textbf{Benchmarks.} We evaluate three tracks. The \emph{GIS track} has 125 questions: 85 Spatial questions across three difficulty levels and 40 Robustness questions spanning security rejection, anti-illusion, schema hallucination, schema enforcement, data tampering prevention, and OOM prevention. The \emph{BIRD track} uses BIRD mini\_dev~\citep{li2023bird} with a 108-question design set and a disjoint 150-question held-out set. The \emph{cross-lingual track} has 50 Chinese-translated BIRD questions plus 50 native Chinese GIS questions.

For external-benchmark readiness, we ported FloodSQL-Bench~\citep{liu2025floodsql} to local PostgreSQL/PostGIS. The port contains 443 questions over 10 flood-risk, social-vulnerability, infrastructure, and boundary tables; 18 are reserved as few-shot examples. A stratified 150-question smoke run with \texttt{gemini-2.5-flash} produced identical baseline and full EX (65/150, 43.3\%). We report it only as a harness-readiness audit, not as generalisation evidence.

\begin{table}[t]
\centering
\footnotesize
\caption{CQ-125 Spatial-subset profile. Predicate families may overlap because one question can require more than one spatial operation. The 40-question Robustness subset is reported separately.}
\label{tab:cq-profile}
% Auto-generated by scripts/nl2sql_bench_cq/benchmark_profile.py
% Regenerate after any benchmark change; do not hand-edit.
\begin{tabular}{lr}
\toprule
Category & Count \\
\midrule
\multicolumn{2}{l}{\emph{Predicate family (a query may belong to several)}} \\
\quad topological & 21 \\
\quad metric & 20 \\
\quad knn & 5 \\
\quad aggregation & 46 \\
\quad plain & 23 \\
\multicolumn{2}{l}{\emph{PostGIS function frequency}} \\
\quad \texttt{ST\_Transform} & 14 \\
\quad \texttt{ST\_Intersects} & 12 \\
\quad \texttt{ST\_Area} & 7 \\
\quad \texttt{ST\_Length} & 6 \\
\quad \texttt{ST\_Contains} & 6 \\
\quad \texttt{ST\_Distance} & 5 \\
\quad \texttt{<->} & 5 \\
\quad \texttt{ST\_AsText} & 3 \\
\quad \texttt{ST\_Within} & 3 \\
\quad \texttt{ST\_Centroid} & 2 \\
\quad \texttt{ST\_DWithin} & 2 \\
\quad \texttt{ST\_Union} & 2 \\
\multicolumn{2}{l}{\emph{Multi-table join multiplicity}} \\
\quad single-table & 55 \\
\quad 2-table join & 29 \\
\quad 3+ table join & 1 \\
\multicolumn{2}{l}{\emph{Difficulty stratification}} \\
\quad Easy & 24 \\
\quad Medium & 36 \\
\quad Hard & 25 \\
\midrule
Total Spatial questions & 85 \\
\bottomrule
\end{tabular}

\end{table}

EX is the primary metric: a prediction is correct only when its result set matches the gold SQL under set-equality with numeric tolerance. We also report Wilson intervals, execution validity, difficulty breakdowns, and paired McNemar tests. For repeated stochastic samples in \S\ref{sec:gisr}, the question is the unit of inference: majority vote is taken per question and condition before one exact two-sided McNemar test. Repeated samples are not pooled as independent observations.

\paragraph{Execution accuracy evaluator.}
EX compares predicted and gold result sets with row-order invariance, positional column order, numeric tolerance ($10^{-3}$ absolute or relative), NULL equality, three-decimal float rounding, gold-\texttt{LIMIT} handling, and duplicate-row preservation. The evaluator is a single Python function (\texttt{compare\_results()}) in the companion repository; all 85 Spatial gold queries return at most 10\,000 rows.

For the framework-level tracks, we compare two pipeline configurations:
\begin{itemize}[leftmargin=*]
  \item \textbf{Baseline}: direct LLM generation (Gemini~2.5~Flash) with schema dump only, no semantic grounding.
  \item \textbf{Full pipeline}: full NL2Semantic2SQL with semantic-layer resolution, geometry-aware grounding, value-aware hints, postprocessing, and execution-time self-correction (single-pass on BIRD, agent-loop on GIS).
\end{itemize}

We additionally compare against \textbf{DIN-SQL}~\citep{pourreza2023dinsql} on framework-level tracks. These supporting configurations share the same database backend and evaluator. The cross-family grounding-effect study varies the LLM family while keeping A/B prompt conditions, benchmark, evaluator, and question-level test convention fixed. Full configurations are in Supplement~S1.
